\worksheettitle
  {Точка, прямая, отрезок и луч}
  {\modulemark{M01} Учебный модуль \textbullet\ полный маршрут}

\studentnameblock

\begin{notebox}[Чему вы научитесь]
  Различать точку, прямую, отрезок и луч, правильно называть их и выполнять
  построения по обозначениям.
\end{notebox}

\begin{checkpointbox}[Вход: проверьте опору]
  \begin{enumerate}[label=\textbf{\arabic*.}]
    \item Отметьте две точки и обозначьте их буквами $A$ и $B$.
      \workarea{12mm}
    \item Какой инструмент нужен, чтобы провести через них прямую?
      \answerblank[45mm]
    \item Какой из объектов имеет два конца: прямая, отрезок или луч?
      \answerblank[35mm]
  \end{enumerate}
\end{checkpointbox}

\begin{notebox}[Как устроен маршрут]
  Сначала разберите образец сравнения четырёх объектов. Затем выполните
  задания с рисунками и построениями, исправьте две типичные ошибки и пройдите
  контрольную точку. Дополнительные задания выполняйте по указанию учителя.
\end{notebox}

\section{Разбираем пример}

\objectscomparisonfigure

\begin{fillbox}[Что сделали]
  Смотрим не на длину нарисованной линии, а на её устройство.
  У прямой нет \answerblank[24mm], у отрезка их \answerblank[16mm], у луча
  есть только \answerblank[24mm]. Точка показывает положение и обозначается
  \answerblank[38mm] латинской буквой.
\end{fillbox}

\begin{notebox}[Правило сравнения]
  \textbf{Прямая} продолжается в обе стороны и не имеет концов.
  \textbf{Отрезок} имеет два конца. \textbf{Луч} имеет начало и продолжается
  в одну сторону. В названии луча первая буква обозначает его начало.
\end{notebox}

\begin{center}
  \renewcommand{\arraystretch}{1.35}
  \begin{tabularx}{\textwidth}{>{\raggedright\arraybackslash}p{0.19\textwidth}
    >{\centering\arraybackslash}p{0.19\textwidth}
    >{\raggedright\arraybackslash}p{0.34\textwidth}
    >{\centering\arraybackslash}X}
    \toprule
    \rowcolor{TableHeader}
    Объект & Концы & Продолжение & Можно измерить длину? \\
    \midrule
    Прямая & нет & в обе стороны & нет \\
    Отрезок & два & не продолжается за концы & да \\
    Луч & одно начало & в одну сторону & нет \\
    \bottomrule
  \end{tabularx}
\end{center}

\section{Практика с опорой}

\incidenceexamplefigure

\begin{practicebox}[1. Точки и прямая]
  \begin{enumerate}[label=\alph*)]
    \item На прямой $a$ лежат точки \answerblank[32mm].
    \item Точка \answerblank[15mm] не лежит на прямой $a$.
    \item Прямую можно назвать по двум точкам: прямая \answerblank[22mm].
  \end{enumerate}
\end{practicebox}

\threelinepointsfigure

\begin{practicebox}[2. Все отрезки]
  Назовите все отрезки, концами которых служат отмеченные точки:
  \[
    \answerblank[25mm],\qquad \answerblank[25mm],\qquad
    \answerblank[25mm].
  \]
\end{practicebox}

\oppositeraysfigure

\begin{practicebox}[3. Лучи с общим началом]
  Луч, идущий из $A$ через $B$: \answerblank[24mm].
  Луч, идущий из $A$ через $C$: \answerblank[24mm].
  Их общее начало: \answerblank[20mm].
\end{practicebox}

\section{Самостоятельная практика}

\identifyobjectsfigure

\begin{practicebox}[4. Назовите объекты]
  Запишите под каждым рисунком полное название: например, «отрезок $AB$».
\end{practicebox}

\begin{practicebox}[5. Выполните построения линейкой]
  Проведите прямую $AB$, отрезок $CD$ и луч $DA$.
  \constructionpointsfigure
\end{practicebox}

\section{Разбираем ошибки}

\begin{warningbox}[Ученик записал неверно]
  Ученик сказал: «Луч $AB$ начинается в точке $B$, а отрезки $AB$ и
  $BA$ разные».

  \textbf{Причина ошибки:} \answerblank[100mm]

  \textbf{Исправление:} \writinglines{2}
\end{warningbox}

\section{Контрольная точка}

\begin{checkpointbox}[6. Проверьте способ на новых данных]
  \begin{enumerate}[label=\textbf{\arabic*.}]
    \item Какой объект имеет начало, но не имеет конца? \answerblank[32mm]
    \item Верно ли, что прямая $MN$ заканчивается в точках $M$ и $N$?
      Объясните: \answerblank[80mm]
    \item Луч начинается в $K$ и проходит через $P$. Запишите его название:
      \answerblank[25mm].
    \item На прямой отмечены точки $D,E,F$. Сколько отрезков имеют концы в
      этих точках? Назовите их: \answerblank[75mm].
  \end{enumerate}
\end{checkpointbox}

\begin{reflectionbox}[Проверяю себя]
  Отметьте выполненные действия.
  \begin{itemize}[label=\choicebox]
    \item различаю прямую, отрезок и луч по числу концов;
    \item называю луч, начиная с его начала;
    \item называю все отрезки по отмеченным точкам;
    \item строю объект по его обозначению.
  \end{itemize}
\end{reflectionbox}

\begin{resultbox}[Дальнейший маршрут]
  Если путаете число концов или порядок букв в луче, выполните M01-R.
  Если способ понятен, но нужна тренировка, выполните M01-H. При устойчивом
  результате переходите к M02 или M05. M01\(\star\) выполняется дополнительно.
\end{resultbox}
